%Version 3.1 December 2024
% See section 11 of the User Manual for version history
%
%%%%%%%%%%%%%%%%%%%%%%%%%%%%%%%%%%%%%%%%%%%%%%%%%%%%%%%%%%%%%%%%%%%%%%
%%                                                                 %%
%% Please do not use \input{...} to include other tex files.       %%
%% Submit your LaTeX manuscript as one .tex document.              %%
%%                                                                 %%
%% All additional figures and files should be attached             %%
%% separately and not embedded in the \TeX\ document itself.       %%
%%                                                                 %%
%%%%%%%%%%%%%%%%%%%%%%%%%%%%%%%%%%%%%%%%%%%%%%%%%%%%%%%%%%%%%%%%%%%%%

%%\documentclass[referee,sn-basic]{sn-jnl}% referee option is meant for double line spacing

%%=======================================================%%
%% to print line numbers in the margin use lineno option %%
%%=======================================================%%

%%\documentclass[lineno,pdflatex,sn-basic]{sn-jnl}% Basic Springer Nature Reference Style/Chemistry Reference Style

%%=========================================================================================%%
%% the documentclass is set to pdflatex as default. You can delete it if not appropriate.  %%
%%=========================================================================================%%

%%\documentclass[sn-basic]{sn-jnl}% Basic Springer Nature Reference Style/Chemistry Reference Style

%%Note: the following reference styles support Namedate and Numbered referencing. By default the style follows the most common style. To switch between the options you can add or remove Numbered in the optional parenthesis. 
%%The option is available for: sn-basic.bst, sn-chicago.bst%  
 
%%\documentclass[pdflatex,sn-nature]{sn-jnl}% Style for submissions to Nature Portfolio journals
%%\documentclass[pdflatex,sn-basic]{sn-jnl}% Basic Springer Nature Reference Style/Chemistry Reference Style
\documentclass[pdflatex,sn-mathphys-num]{sn-jnl}% Math and Physical Sciences Numbered Reference Style
%%\documentclass[pdflatex,sn-mathphys-ay]{sn-jnl}% Math and Physical Sciences Author Year Reference Style
%%\documentclass[pdflatex,sn-aps]{sn-jnl}% American Physical Society (APS) Reference Style
%%\documentclass[pdflatex,sn-vancouver-num]{sn-jnl}% Vancouver Numbered Reference Style
%%\documentclass[pdflatex,sn-vancouver-ay]{sn-jnl}% Vancouver Author Year Reference Style
%%\documentclass[pdflatex,sn-apa]{sn-jnl}% APA Reference Style
%%\documentclass[pdflatex,sn-chicago]{sn-jnl}% Chicago-based Humanities Reference Style

%%%% Standard Packages
%%<additional latex packages if required can be included here>

\usepackage{graphicx}%
\usepackage{multirow}%
\usepackage{amsmath,amssymb,amsfonts}%
\usepackage{amsthm}%
\usepackage{mathrsfs}%
\usepackage[title]{appendix}%
\usepackage{xcolor}%
\usepackage{textcomp}%
\usepackage{manyfoot}%
\usepackage{booktabs}%
\usepackage{algorithm}%
\usepackage{algorithmicx}%
\usepackage{algpseudocode}%
\usepackage{listings}%

%%%%

%%%%%=============================================================================%%%%
%%%%  Remarks: This template is provided to aid authors with the preparation
%%%%  of original research articles intended for submission to journals published 
%%%%  by Springer Nature. The guidance has been prepared in partnership with 
%%%%  production teams to conform to Springer Nature technical requirements. 
%%%%  Editorial and presentation requirements differ among journal portfolios and 
%%%%  research disciplines. You may find sections in this template are irrelevant 
%%%%  to your work and are empowered to omit any such section if allowed by the 
%%%%  journal you intend to submit to. The submission guidelines and policies 
%%%%  of the journal take precedence. A detailed User Manual is available in the 
%%%%  template package for technical guidance.
%%%%%=============================================================================%%%%

%% as per the requirement new theorem styles can be included as shown below
\theoremstyle{thmstyleone}%
\newtheorem{theorem}{Theorem}%  meant for continuous numbers
%%\newtheorem{theorem}{Theorem}[section]% meant for sectionwise numbers
%% optional argument [theorem] produces theorem numbering sequence instead of independent numbers for Proposition
\newtheorem{proposition}[theorem]{Proposition}% 
%%\newtheorem{proposition}{Proposition}% to get separate numbers for theorem and proposition etc.

\theoremstyle{thmstyletwo}%
\newtheorem{example}{Example}%
\newtheorem{remark}{Remark}%

\theoremstyle{thmstylethree}%
\newtheorem{definition}{Definition}%

\raggedbottom
%%\unnumbered% uncomment this for unnumbered level heads

\begin{document}

\title[Article Title]{An Explainable Multi-Class Disaster Tweet Classification System with Geospatial Temporal Aggregation and Human-in-the-Loop Verification}

%%=============================================================%%
%% GivenName	-> \fnm{Joergen W.}
%% Particle	-> \spfx{van der} -> surname prefix
%% FamilyName	-> \sur{Ploeg}
%% Suffix	-> \sfx{IV}
%% \author*[1,2]{\fnm{Joergen W.} \spfx{van der} \sur{Ploeg} 
%%  \sfx{IV}}\email{iauthor@gmail.com}
%%=============================================================%%

\author*[1]{\fnm{Vikram Dharsan} \sur{K S}}
\email{2023103050@student.annauniv.edu}

\author*[1]{\fnm{Abhijith} \sur{M}}
\email{2023103095@student.annauniv.edu}

\author*[1]{\fnm{Vijay Anandan} \sur{J M}}
\email{2023103541@student.annauniv.edu}

\author*[1]{\fnm{Jeeva} \sur{S}}
\email{2023103604@student.annauniv.edu}

\affil*[1]{\orgdiv{Department of Computer Science Engineering},
\orgname{College of Engineering, Guindy},
\orgaddress{
\city{Chennai},
\country{India}
}}

%%==================================%%
%% Sample for unstructured abstract %%
%%==================================%%

\abstract{Social media platforms such as Twitter have become critical real-time information sources during natural disasters; however, the large volume of tweets makes manual analysis infeasible. Existing approaches primarily rely on binary classification, which lacks the granularity required for effective humanitarian response. This paper presents an end-to-end disaster tweet classification system that addresses three key challenges: multi-class categorization, model interpretability, and location-aware situational analysis. The system classifies tweets into five humanitarian categories using a fine-tuned Sentence-BERT (all-MiniLM-L6-v2) model with a lightweight classification head. It integrates gradient-based Explainable AI (XAI) for token-level interpretation, a spaCy-based Named Entity Recognition (NER) module for extracting actionable information, and a Human-in-the-Loop (HITL) mechanism for low-confidence predictions. Additionally, a Geospatial Temporal Aggregation module clusters tweets by location, computes consensus distributions, performs temporal analysis, and generates situation reports with severity scoring and recommended actions. Experimental results achieve an accuracy of 87.0\% and a macro F1-score of 87.26\% on multi-event disaster data. The system operates entirely on local hardware, ensuring offline capability.  }

%%================================%%
%% Sample for structured abstract %%
%%================================%%

% \abstract{\textbf{Purpose:} The abstract serves both as a general introduction to the topic and as a brief, non-technical summary of the main results and their implications. The abstract must not include subheadings (unless expressly permitted in the journal's Instructions to Authors), equations or citations. As a guide the abstract should not exceed 200 words. Most journals do not set a hard limit however authors are advised to check the author instructions for the journal they are submitting to.
% 
% \textbf{Methods:} The abstract serves both as a general introduction to the topic and as a brief, non-technical summary of the main results and their implications. The abstract must not include subheadings (unless expressly permitted in the journal's Instructions to Authors), equations or citations. As a guide the abstract should not exceed 200 words. Most journals do not set a hard limit however authors are advised to check the author instructions for the journal they are submitting to.
% 
% \textbf{Results:} The abstract serves both as a general introduction to the topic and as a brief, non-technical summary of the main results and their implications. The abstract must not include subheadings (unless expressly permitted in the journal's Instructions to Authors), equations or citations. As a guide the abstract should not exceed 200 words. Most journals do not set a hard limit however authors are advised to check the author instructions for the journal they are submitting to.
% 
% \textbf{Conclusion:} The abstract serves both as a general introduction to the topic and as a brief, non-technical summary of the main results and their implications. The abstract must not include subheadings (unless expressly permitted in the journal's Instructions to Authors), equations or citations. As a guide the abstract should not exceed 200 words. Most journals do not set a hard limit however authors are advised to check the author instructions for the journal they are submitting to.}

\keywords{Disaster Tweet Classification, Multi-Class Classification, Explainable AI, Geospatial Aggregation, Transformer Models, Human-in-the-Loop, Named Entity Recognition, Crisis Informatics, Natural Language Processing }

%%\pacs[JEL Classification]{D8, H51}

%%\pacs[MSC Classification]{35A01, 65L10, 65L12, 65L20, 65L70}

\maketitle

\section{Introduction}

Natural disasters such as earthquakes, hurricanes, floods, and wildfires generate an unprecedented volume of social media activity. During Hurricane Harvey (2017), over 325,000 disaster-related tweets were posted within the first 48 hours \cite{ref1}. These tweets contain critical, time-sensitive information, including reports of trapped individuals, infrastructure damage, resource needs, and rescue coordination. For emergency responders, government agencies, and humanitarian organizations, social media has become an indispensable source of real-time situational awareness \cite{ref2}.

However, the sheer volume of tweets during a crisis event renders manual analysis infeasible. Traditional keyword-based filtering produces high false-positive rates, while binary classification systems (disaster vs.\ non-disaster) fail to capture the semantic diversity of crisis-related information \cite{ref3}. Emergency responders need to know not just whether a tweet is disaster-related, but specifically what type of information it contains, such as casualty reports, infrastructure damage, or rescue coordination.

\subsection{Problem Statement}

Existing disaster tweet classification systems suffer from three primary limitations:

\begin{enumerate}
    \item \textbf{Insufficient Granularity:} Binary classification cannot differentiate between tweets reporting casualties, infrastructure damage, and rescue efforts, each requiring different response actions.
    
    \item \textbf{Lack of Interpretability:} Deep learning models often provide predictions without explanations, making it difficult for operators to trust and validate automated decisions.
    
    \item \textbf{Absence of Spatial Aggregation:} Tweet-level predictions are not aggregated into location-level insights, limiting their usefulness for coordinated response.
\end{enumerate}

\subsection{Contributions}

This paper makes the following contributions:

\begin{itemize}
    \item \textbf{End-to-End Classification System:} An integrated system built upon the DeLTran15 \cite{ref16} transformer model that classifies disaster tweets into five humanitarian categories, achieving 87.0\% accuracy and 87.26\% macro F1-score.
    
    \item \textbf{Explainable AI (XAI):} A gradient-based token importance mechanism that provides interpretable predictions through token-level scoring.
    
    \item \textbf{Actionable Information Extraction:} A hybrid NLP pipeline combining spaCy Named Entity Recognition and regex-based methods to extract structured information such as locations, casualty counts, and resource needs.
    
    \item \textbf{Human-in-the-Loop (HITL):} A confidence-based verification workflow that enables human validation for low-confidence predictions.
    
    \item \textbf{Geospatial Temporal Aggregation:} A multi-step pipeline that clusters tweets by location, performs temporal analysis, and generates situation reports with severity scoring and recommended actions.
\end{itemize}
\subsection{Paper Organization}

The remainder of this paper is organized as follows. Section II reviews related work in disaster tweet classification, explainable AI, geospatial analysis, and human-in-the-loop systems. Section III describes the dataset and preprocessing methodology used in this study. Section IV presents the system architecture built upon DeLTran15, including the classification model, explainable AI module, actionable information extraction pipeline, and the geospatial temporal aggregation module with its location resolution, consensus computation, severity scoring, and report generation components. Section V details the experimental setup, training configuration, and evaluation results. Section VI describes the system implementation, including the client-server architecture and deployment details. Section VII provides a comparative study with existing methods along with performance visualizations and key insights. Finally, Section VIII concludes the paper and outlines directions for future work.

\section{Related Work}

\subsection{Disaster Tweet Classification}

Early approaches to disaster tweet classification relied on traditional machine learning methods. Imran et al. \cite{ref4} employed Support Vector Machines (SVM) with TF-IDF features to classify crisis tweets, achieving F1-scores between 65--75\% on crisis datasets. Alam et al. \cite{ref5} introduced the CrisisMMD dataset, a multi-modal benchmark containing tweets from multiple disaster events annotated with humanitarian categories.

The advent of transformer-based models significantly improved classification performance. Devlin et al. \cite{ref6} demonstrated that BERT-based models achieve state-of-the-art results on text classification benchmarks. Liu et al. \cite{ref7} further showed that domain-specific fine-tuning of pre-trained language models yields substantial improvements over traditional feature engineering approaches.

\subsection{Explainable AI in NLP}

The need for model interpretability in safety-critical applications has driven research in Explainable AI (XAI). Ribeiro et al. \cite{ref8} proposed LIME (Local Interpretable Model-agnostic Explanations), which approximates model behavior locally using interpretable surrogate models. Lundberg and Lee \cite{ref9} introduced SHAP (SHapley Additive exPlanations), based on cooperative game theory. Gradient-based methods, including Integrated Gradients \cite{ref10} and attention-based explanations \cite{ref11}, provide token-level importance scores directly from model representations without requiring surrogate models.

\subsection{Geospatial Analysis in Crisis Informatics}

Geospatial analysis of social media during disasters has been explored in prior work. De Albuquerque et al. \cite{ref12} demonstrated the correlation between tweet density and flood impact areas. Middleton et al. \cite{ref13} developed real-time geoparsing systems for crisis events. However, existing approaches primarily focus on binary disaster detection per location and do not provide multi-class consensus-based situational awareness, which is addressed in our work.

\subsection{Human-in-the-Loop Systems}

Human-in-the-Loop (HITL) approaches have been shown to improve model reliability in critical applications. Branson et al. \cite{ref14} demonstrated that human verification enhances classification performance. Monarch \cite{ref15} formalized HITL workflows for iterative model improvement. In this work, we adopt a confidence-threshold-based HITL mechanism, where predictions below a predefined threshold are flagged for human review.

\section{Dataset and Preprocessing}

\subsection{Dataset Description}

We utilize the CrisisMMD (Crisis Multi-Modal Dataset) \cite{ref5}, comprising tweets from seven major disaster events, as shown in Table~\ref{tab:dataset}.

\begin{table}[h]
\centering
\caption{CrisisMMD Dataset Overview}
\label{tab:dataset}
\begin{tabular}{|l|c|l|c|}
\hline
\textbf{Disaster Event} & \textbf{Year} & \textbf{Region} & \textbf{Tweet Count} \\
\hline
Hurricane Harvey & 2017 & Texas, USA & $\sim$5,200 \\
Hurricane Irma & 2017 & Caribbean/Florida & $\sim$5,300 \\
Hurricane Maria & 2017 & Puerto Rico & $\sim$5,400 \\
California Wildfires & 2017 & California, USA & $\sim$1,900 \\
Mexico Earthquake & 2017 & Mexico & $\sim$1,700 \\
Iraq-Iran Earthquake & 2017 & Iraq/Iran border & $\sim$700 \\
Sri Lanka Floods & 2017 & Sri Lanka & $\sim$1,100 \\
\hline
\end{tabular}
\end{table}

Each tweet is annotated with one of five humanitarian categories, as shown in Table~\ref{tab:labels}.

\begin{table}[h]
\centering
\caption{Humanitarian Categories}
\label{tab:labels}
\begin{tabular}{|c|l|l|}
\hline
\textbf{Label ID} & \textbf{Category} & \textbf{Description} \\
\hline
0 & Affected Individuals & Casualties, missing persons, displacement \\
1 & Infrastructure \& Utility Damage & Damage to roads, bridges, power, water \\
2 & Not Humanitarian & Irrelevant or non-disaster content \\
3 & Other Relevant Information & Warnings, updates, awareness \\
4 & Rescue/Volunteering/Donation & Help offers, shelters, donations \\
\hline
\end{tabular}
\end{table}

\subsection{Data Preprocessing}

The preprocessing pipeline applies the following transformations:

\begin{enumerate}
    \item \textbf{URL Removal:} All HTTP/HTTPS URLs are removed as they do not contribute to semantic classification.
    
    \item \textbf{Special Character Handling:} @mentions are removed, and hashtag symbols are stripped while preserving the textual content (e.g., \#flood $\rightarrow$ flood).
    
    \item \textbf{Contraction Expansion:} English contractions are expanded (e.g., ``don't'' $\rightarrow$ ``do not'') to normalize vocabulary.
    
    \item \textbf{Emoji Translation:} Emojis are converted to textual descriptions to preserve sentiment information.
    
    \item \textbf{Whitespace Normalization:} Multiple spaces are reduced to single spaces, and leading/trailing spaces are removed.
    
    \item \textbf{Case Normalization:} All text is converted to lowercase.
\end{enumerate}

The preprocessing pipeline reduces vocabulary noise while preserving semantic information. The cleaned dataset is stored as \texttt{disaster\_text\_clean.csv}, containing processed text and corresponding labels.

\section{System Architecture}

\subsection{Overview}

The system follows a modular, privacy-first architecture where all AI inference runs locally on the user's machine. It is built upon the DeLTran15 classification model \cite{ref16} and extends it with five additional modules: (i) the Explainable AI engine, (ii) the Actionable Information Extraction module, (iii) the Human-in-the-Loop verification workflow, (iv) the Geospatial Temporal Aggregation pipeline, and (v) the Severity Scoring and Action Recommendation engine.

The overall architecture of the system is illustrated in Fig.~\ref{fig:architecture}.

\begin{figure*}[t]
    \centering
    \includegraphics[width=\textwidth]{CIP2.jpeg}
    \caption{System architecture illustrating preprocessing, DeLTran15-based classification, explainable AI, actionable information extraction, and geospatial temporal aggregation.}
    \label{fig:architecture}
\end{figure*}

\subsection{DeLTran15 Classification Model}

The classification backbone adopts the DeLTran15 model \cite{ref16}, which is built upon the \texttt{sentence-transformers/all-MiniLM-L6-v2} pre-trained transformer, selected for its strong performance on short-text semantic understanding tasks while maintaining a lightweight footprint.

The architecture consists of:

\begin{itemize}
    \item \textbf{Transformer Encoder:} A 6-layer transformer with 384-dimensional hidden representations.
    \item \textbf{Regularization Layer:} Dropout with probability $p = 0.3$ applied to the [CLS] embedding.
    \item \textbf{Classification Head:} A linear layer mapping the 384-dimensional embedding to 5 output logits.
\end{itemize}

Formally, for an input tweet $t$:

\begin{equation}
h = \text{Encoder}(\text{Tokenize}(t))
\end{equation}

\begin{equation}
c = h_{[CLS]}
\end{equation}

\begin{equation}
z = \text{Dropout}(c, p=0.3)
\end{equation}

\begin{equation}
y = \text{Softmax}(Wz + b)
\end{equation}

where $W \in \mathbb{R}^{5 \times 384}$ and $b \in \mathbb{R}^{5}$ are learnable parameters.

\subsection{Explainable AI (XAI) Module}

To provide interpretability, we employ gradient-based token importance analysis. After the forward pass, gradients of the predicted class are computed with respect to input embeddings. The importance of each token $i$ is defined as:

\begin{equation}
\text{importance}(i) = \|h_i\|_2 = \sqrt{\sum_j h_{ij}^2}
\end{equation}

where $h_i \in \mathbb{R}^{384}$ is the hidden state vector for token $i$. The scores are normalized to $[0,1]$ and visualized using a color gradient.

\subsection{Actionable Information Extraction}

For disaster-related tweets (labels 0, 1, and 4), a hybrid NLP pipeline extracts actionable intelligence:

\begin{enumerate}
    \item \textbf{Location Extraction:} spaCy \texttt{en\_core\_web\_sm} model identifies GPE, LOC, and FAC entities, with regex-based fallback patterns.
    
    \item \textbf{Casualty Count Extraction:} Regex patterns detect phrases such as ``30 injured'' or ``12 dead'' using predefined status keywords.
    
    \item \textbf{Resource Need Identification:} Keyword matching is performed over a predefined vocabulary (e.g., food, water, shelter, medical aid).
    
    \item \textbf{Damage Type Classification:} Keywords such as bridge, road, power, flooded, and collapsed are used to detect infrastructure damage.
    
    \item \textbf{Temporal Reference Detection:} Regex patterns identify time indicators such as ``now'', ``today'', or ``N hours ago''.
\end{enumerate}

\subsection{Human-in-the-Loop Verification}

The HITL module implements a confidence-based verification workflow:

\begin{itemize}
    \item If $\max(\text{Softmax}) \geq 0.70$, the prediction is automatically accepted as \textit{verified}.
    
    \item If $\max(\text{Softmax}) < 0.70$, the prediction is marked as \textit{unverified} and queued for human review.
\end{itemize}

The reviewer is provided with the tweet text, predicted label, confidence score, and XAI-based token highlighting. Corrections are stored as \textit{human\_verified} and can be used for future model improvement.

\subsection{Geospatial Temporal Aggregation}

Individual tweet-level classification is insufficient for effective disaster response coordination. Emergency managers require location-level situational awareness, such as understanding the overall condition of a region and identifying whether events are active, recovering, or unreliable. This module addresses this need by clustering tweets by geographic location and generating per-location situation reports.

The overall geospatial temporal aggregation pipeline is illustrated in Fig.~\ref{fig:geo_pipeline}.

\begin{figure*}[t]
\centering
\includegraphics[width=\textwidth]{GSA.png}
\caption{Geospatial temporal aggregation pipeline showing classification, location resolution, clustering, temporal analysis, decision matrix, severity scoring, and report generation.}
\label{fig:geo_pipeline}
\end{figure*}

\subsubsection{Location Resolution}

Each tweet's location is resolved using a three-tier priority system:

\begin{enumerate}
    \item \textbf{Place Tag:} GPS-based geolocation (most reliable, limited availability).
    \item \textbf{User Profile Location:} Self-declared user location.
    \item \textbf{Text Extraction:} Locations extracted using spaCy Named Entity Recognition.
\end{enumerate}

Location strings are normalized by converting to lowercase and removing prefixes/suffixes to ensure consistent clustering.

\subsubsection{Six-Step Aggregation Pipeline}

For each location cluster containing $N$ tweets, the following steps are executed:

\textbf{Step 1 — Cluster Formation:} Tweets are grouped by normalized location.

\textbf{Step 2 — Consensus Computation:}

\begin{equation}
\text{consensus}(k) = \frac{\sum_i (\text{confidence}_i \cdot \mathbb{1}[y_i = k])}{\sum_i \text{confidence}_i}
\end{equation}

This produces a distribution across the five humanitarian categories.

\textbf{Step 3 — Temporal Analysis:}

\begin{itemize}
    \item Burst ($\leq$ 6 hours)
    \item Spread ($>$ 72 hours)
    \item Normal (6–72 hours)
\end{itemize}

\textbf{Step 4 — Decision Classification:}

\begin{table}[h]
\centering
\caption{Decision Matrix for Cluster Classification}
\label{tab:decision}
\begin{tabular}{|c|l|l|l|}
\hline
\textbf{Priority} & \textbf{Case} & \textbf{Condition} & \textbf{Severity} \\
\hline
1 & Insufficient Data & Single tweet & UNKNOWN \\
2 & Suspicious Source & Single author & UNKNOWN \\
3 & Early Signal & $<$ 5 tweets & LOW \\
4 & Location Uncertain & Mixed locations & UNCERTAIN \\
5 & No Disaster & $\geq$ 80\% non-disaster & NONE \\
6 & Low Confidence & Avg $<$ 50\% & UNCERTAIN \\
7 & Ambiguous & Low agreement & MEDIUM \\
8 & Active Event & Burst + high disaster & CRITICAL \\
9 & Recovery Phase & Spread + recovery & MEDIUM \\
10 & Confirmed Event & Meets thresholds & COMPUTED \\
\hline
\end{tabular}
\end{table}

\textbf{Step 5 — Intelligence Merging:} Extracted entities such as locations, needs, damage types, and people counts are merged and deduplicated.

\textbf{Step 6 — Action Recommendation:} Based on dominant labels:
\begin{itemize}
    \item Affected Individuals $\rightarrow$ Deploy rescue teams
    \item Infrastructure Damage $\rightarrow$ Send repair teams
    \item Rescue/Donation $\rightarrow$ Coordinate aid distribution
\end{itemize}

\subsubsection{Severity Scoring}

For confirmed events, a point-based severity scoring mechanism is applied:

\begin{equation}
\text{Severity Score} = f(\text{tweet count}, \text{people affected}, \text{damage reports})
\end{equation}

The final severity is categorized as:
\begin{itemize}
    \item CRITICAL
    \item HIGH
    \item MEDIUM
    \item LOW
\end{itemize}

\section{Experimental Results}

\subsection{Training Configuration}

The model was fine-tuned using the hyperparameters listed in Table~\ref{tab:train_config}.

\begin{table}[h]
\centering
\caption{Training Configuration}
\label{tab:train_config}
\begin{tabular}{|l|l|}
\hline
\textbf{Parameter} & \textbf{Value} \\
\hline
Base Model & all-MiniLM-L6-v2 \\
Optimizer & AdamW \\
Learning Rate & $2 \times 10^{-5}$ \\
Batch Size & 32 \\
Epochs & 15 \\
Dropout Rate & 0.3 \\
Max Sequence Length & 128 tokens \\
Loss Function & Cross-Entropy \\
\hline
\end{tabular}
\end{table}

\subsection{Classification Performance}

The model was evaluated on a held-out test set of 300 labeled tweets across multiple disaster events. The overall performance is summarized in Table~\ref{tab:overall_metrics}.

\begin{table}[h]
\centering
\caption{Overall Performance Metrics}
\label{tab:overall_metrics}
\begin{tabular}{|l|c|}
\hline
\textbf{Metric} & \textbf{Score} \\
\hline
Accuracy & 87.0\% \\
Macro F1-Score & 87.26\% \\
Weighted F1-Score & 87.26\% \\
Macro Precision & 88.15\% \\
Macro Recall & 86.92\% \\
\hline
\end{tabular}
\end{table}

The per-class performance is shown in Table~\ref{tab:class_metrics}.

\begin{table}[h]
\centering
\caption{Per-Class Performance}
\label{tab:class_metrics}
\begin{tabular}{|l|c|c|c|c|}
\hline
\textbf{Class} & \textbf{Precision} & \textbf{Recall} & \textbf{F1} & \textbf{Support} \\
\hline
Affected Individuals & 0.85 & 0.83 & 0.84 & 58 \\
Infrastructure Damage & 0.90 & 0.88 & 0.89 & 62 \\
Not Humanitarian & 0.92 & 0.94 & 0.93 & 65 \\
Other Information & 0.83 & 0.85 & 0.84 & 55 \\
Rescue/Donation & 0.91 & 0.85 & 0.88 & 60 \\
\hline
\end{tabular}
\end{table}

\subsection{Confidence Distribution Analysis}

Analysis of the model's confidence scores shows a mean confidence of 0.89 and a median of 0.93. Approximately 78\% of predictions exceed the 0.70 threshold used in the Human-in-the-Loop (HITL) module, allowing automatic acceptance. The remaining 22\% are flagged for human review, corresponding to ambiguous or borderline cases.

\subsection{Actionable Information Extraction Coverage}

For disaster-related tweets, the NLP extraction module achieves the coverage shown in Table~\ref{tab:coverage}.

\begin{table}[h]
\centering
\caption{Actionable Information Extraction Coverage}
\label{tab:coverage}
\begin{tabular}{|l|c|}
\hline
\textbf{Field} & \textbf{Coverage Rate} \\
\hline
Location Extraction & 72\% \\
People Count & 34\% \\
Resource Needs & 48\% \\
Damage Type & 55\% \\
Time Mentions & 41\% \\
\hline
\end{tabular}
\end{table}

\section{System Implementation}

\subsection{Overview}

The system follows a two-tier architecture comprising a Python-based desktop client and a Node.js backend server. The overall implementation workflow is illustrated in Fig.~\ref{fig:system_impl}.

\begin{figure*}[t]
\centering
\includegraphics[width=\textwidth]{bgcip.png}
\caption{High-level system implementation architecture showing the interaction between the desktop dashboard, DeLTran15 model components, geospatial aggregation module, and backend server with MongoDB storage.}
\label{fig:system_impl}
\end{figure*}

\subsection{Technology Stack}

The system is implemented using PyTorch and HuggingFace Transformers for model inference, spaCy for NLP extraction, CustomTkinter for the desktop interface, Node.js with Express for backend services, and MongoDB for data storage. Communication between components is handled via REST APIs.

\subsection{Privacy-First Design}

All model inference is executed locally on the user's machine. The model weights are stored locally, and no tweet data is transmitted to external services, ensuring data privacy and offline capability.

\subsection{Client-Server Architecture}

The backend provides RESTful endpoints for tweet management and classification updates, while the frontend desktop application interacts with the server using HTTP requests. This decoupled design enables independent deployment and future extensibility.

\subsection{End-to-End Pipeline}

Upon receiving a tweet, the system performs preprocessing, tokenization, and transformer-based classification. Confidence scores are computed using softmax, followed by gradient-based XAI for token importance and NLP-based extraction of actionable information. Predictions below a confidence threshold are flagged for human verification. The final results are stored in MongoDB via API calls.

\subsection{Human-in-the-Loop Feedback}

User corrections are stored in the database and additionally logged in a local CSV file for future retraining, enabling continuous model improvement.

\subsection{XAI Visualization}

Token importance scores are normalized and mapped to a white-to-red color gradient for visualization in the UI, enhancing interpretability of model predictions.


\section{Comparative Study}

\subsection{Performance Comparison with Existing Methods}

Table~\ref{tab:comparison} presents a comparison of the adopted DeLTran15 model with existing approaches for disaster tweet classification.

Early traditional machine learning models such as Na\"{i}ve Bayes and SVM rely on TF-IDF features and achieve limited performance. Shallow embedding methods like FastText improve semantic representation, while deep learning models such as CNN, LSTM, and Bi-LSTM further enhance contextual understanding.

Transformer-based models including DistilBERT, BERT, and RoBERTa significantly outperform earlier methods due to contextual embeddings. However, most approaches lack interpretability and human-in-the-loop validation.

The DeLTran15-based system achieves competitive performance while incorporating explainability, HITL validation, and geospatial analysis.

\begin{table*}[t]
\centering
\caption{Comparative Analysis of Disaster Tweet Classification Methods}
\label{tab:comparison}
\begin{tabular}{|c|l|c|l|c|c|}
\hline
\textbf{\#} & \textbf{Model} & \textbf{Year} & \textbf{Approach} & \textbf{Accuracy} & \textbf{Macro F1} \\
\hline
1 & Na\"{i}ve Bayes & 2017 & Traditional ML & $\sim$0.62 & 0.59 \\
2 & SVM (TF-IDF) & 2018 & Traditional ML & $\sim$0.66 & 0.63 \\
3 & FastText & 2019 & Embedding & 0.76 & 0.794 \\
4 & CNN & 2020 & Deep Learning & 0.80 & 0.842 \\
5 & LSTM & 2020 & Deep Learning & 0.78 & 0.803 \\
6 & Bi-LSTM & 2021 & Deep Learning & 0.80 & 0.820 \\
7 & IDBO CNN-BiLSTM & 2025 & Hybrid DL & 0.803 & 0.798 \\
8 & DistilBERT & 2021 & Transformer & 0.84 & 0.863 \\
9 & BERT (base) & 2021 & Transformer & 0.87 & 0.879 \\
10 & BERTweet & 2022 & Tweet Transformer & 0.85 & 0.855 \\
11 & RoBERTa & 2021 & Transformer & 0.87 & 0.870 \\
12 & CrisisTransformers & 2025 & Domain Transformer & 0.861 & 0.847 \\
13 & He \& Hu Framework & 2025 & BERT + Geo & -- & 0.814 \\
14 & \textbf{DeLTran15 (Ours)} & 2025 & XAI + HITL + Geo & \textbf{0.87} & \textbf{0.8726} \\
\hline
\end{tabular}
\end{table*}

\subsection{Experimental Visualizations}

The model performance and extraction capabilities are further illustrated using the following visualizations.

\begin{figure}[h]
\centering
\includegraphics[width=0.7\linewidth]{fig_predicted_class_distribution.png}
\caption{Predicted class distribution across all test tweets.}
\label{fig:class_dist}
\end{figure}

\begin{figure}[h]
\centering
\includegraphics[width=0.7\linewidth]{fig_actionable_field_coverage.png}
\caption{Coverage of extracted actionable information fields.}
\label{fig:coverage}
\end{figure}

\begin{figure}[h]
\centering
\includegraphics[width=0.7\linewidth]{fig_confusion_matrix_labeled.png}
\caption{Confusion matrix showing classification performance across all categories.}
\label{fig:confusion}
\end{figure}
\subsection{Key Insights}

The adopted DeLTran15 model achieves performance comparable to state-of-the-art transformer models such as RoBERTa while maintaining a significantly smaller model size. With approximately 22M parameters compared to over 300M parameters in larger models, the system achieves substantial efficiency gains.

Additionally, the integration of Explainable AI (XAI), Human-in-the-Loop (HITL) validation, and geospatial aggregation provides enhanced interpretability and real-world applicability, making the system suitable for deployment in disaster response scenarios.



\section{Conclusion and Future Work}

\subsection{Conclusion}

This paper presented a comprehensive disaster tweet classification system built upon the DeLTran15 model \cite{ref16} that addresses the limitations of existing approaches through multi-class categorization, explainable predictions, actionable intelligence extraction, human verification, and geospatial temporal aggregation. The system achieves \textbf{87.0\% accuracy} and \textbf{87.26\% macro F1-score} across five humanitarian categories, while providing location-level situation reports through a \textit{10-case decision matrix}.

The privacy-first architecture ensures that all AI processing runs locally, making the system suitable for deployment in sensitive disaster response environments. The Human-in-the-Loop workflow provides quality assurance for critical predictions while generating labeled data for continuous model improvement.

\subsection{Future Work}

Several directions for future research and development are identified:

\begin{enumerate}
    \item \textbf{Model Enhancement:} Exploring larger transformer architectures (e.g., DeBERTa, XtremeDistil-L12) for improved classification accuracy, particularly for the ``Affected Individuals'' and ``Other Information'' categories which currently show lower F1-scores.

    \item \textbf{Temporal Modeling:} Incorporating sequential tweet analysis to detect evolving disaster narratives over time, enabling real-time trend detection.

    \item \textbf{Multi-modal Classification:} Extending the system to analyze images and videos attached to tweets, providing richer situational awareness.

    \item \textbf{Real-time Streaming:} Integrating with the Twitter Streaming API for continuous, real-time tweet classification and geospatial updating.

    \item \textbf{Cross-lingual Capability:} Adapting the model for disaster tweets in multiple languages using multilingual transformer models (e.g., mBERT, XLM-R).
\end{enumerate}


\begin{thebibliography}{16}

\bibitem{ref1}
M. Imran, C. Castillo, F. Diaz, and S. Vieweg, "Processing social media messages in mass emergency: A survey," \textit{ACM Computing Surveys}, vol. 47, no. 4, pp. 1--38, 2015.

\bibitem{ref2}
A. Olteanu, C. Castillo, F. Diaz, and S. Vieweg, "CrisisLex: A Lexicon for Collecting and Filtering Microblogged Communications in Crises," in \textit{Proc. AAAI Int. Conf. Weblogs Social Media (ICWSM)}, 2014.

\bibitem{ref3}
K. M. Stowe, M. Paul, M. Palmer, L. Palen, and K. Anderson, "Identifying and Categorizing Disaster-Related Tweets," in \textit{Proc. 4th Int. Workshop on Natural Language Processing for Social Media}, 2016, pp. 1--6.

\bibitem{ref4}
M. Imran, S. Elbassuoni, C. Castillo, F. Diaz, and P. Meier, "Extracting information nuggets from disaster-related messages in social media," in \textit{Proc. 10th Int. Conf. Information Systems for Crisis Response and Management (ISCRAM)}, 2013.

\bibitem{ref5}
F. Alam, F. Ofli, and M. Imran, "CrisisMMD: Multimodal Twitter Datasets from Natural Disasters," in \textit{Proc. 12th Int. Conf. on Web and Social Media (ICWSM)}, 2018.

\bibitem{ref6}
J. Devlin, M.-W. Chang, K. Lee, and K. Toutanova, "BERT: Pre-training of Deep Bidirectional Transformers for Language Understanding," in \textit{Proc. NAACL-HLT}, 2019, pp. 4171--4186.

\bibitem{ref7}
Y. Liu, M. Ott, N. Goyal, et al., "RoBERTa: A Robustly Optimized BERT Pretraining Approach," \textit{arXiv preprint arXiv:1907.11692}, 2019.

\bibitem{ref8}
M. T. Ribeiro, S. Singh, and C. Guestrin, "'Why Should I Trust You?': Explaining the Predictions of Any Classifier," in \textit{Proc. ACM SIGKDD}, 2016, pp. 1135--1144.

\bibitem{ref9}
S. M. Lundberg and S.-I. Lee, "A Unified Approach to Interpreting Model Predictions," in \textit{Advances in Neural Information Processing Systems (NeurIPS)}, vol. 30, 2017.

\bibitem{ref10}
M. Sundararajan, A. Taly, and Q. Yan, "Axiomatic Attribution for Deep Networks," in \textit{Proc. ICML}, 2017, pp. 3319--3328.

\bibitem{ref11}
S. Jain and B. C. Wallace, "Attention is not Explanation," in \textit{Proc. NAACL-HLT}, 2019, pp. 3543--3556.

\bibitem{ref12}
J. P. De Albuquerque, B. Herfort, A. Brenning, and A. Zipf, "A geographic approach for combining social media and authoritative data towards identifying useful information for disaster management," \textit{Int. Journal of Geographical Information Science}, vol. 29, no. 4, pp. 667--689, 2015.

\bibitem{ref13}
S. E. Middleton, L. Middleton, and S. Modafferi, "Real-Time Crisis Mapping of Natural Disasters Using Social Media," \textit{IEEE Intelligent Systems}, vol. 29, no. 2, pp. 9--17, 2014.

\bibitem{ref14}
S. Branson, C. Wah, F. Schroff, et al., "Visual Recognition with Humans in the Loop," in \textit{Proc. ECCV}, 2010, pp. 438--451.

\bibitem{ref15}
R. Monarch, \textit{Human-in-the-Loop Machine Learning: Active Learning and Annotation for Human-Centered AI}, Manning Publications, 2021.

\bibitem{ref16}
S. Saleem, N. Hasan, A. Khattar, P. R. Jain, T. K. Gupta, and M. Mehrotra, 
"DeLTran15: A Deep Lightweight Transformer-Based Framework for Multiclass Classification of Disaster Posts on X," 
 

\end{thebibliography}

\end{document}
